\subsection{Fibras formales de los anillos locales noetherianos.}
\label{subsection:IV.7.3-fr}

\begin{env}[7.3.1]
\label{IV.2.s7.3.1-fr}
Vamos a considerar en este número y en los dos siguientes propiedades
$\boldsymbol P(Z,k)$ de la forma siguiente:
\begin{quote}
«$Z$ es un preesquema localmente noetheriano sobre un cuerpo $k$,

y para todo $z\in Z$ se cumple $\boldsymbol Q(\mathscr O_z,k)$».
\end{quote}
Aquí $\boldsymbol Q(A,k)$ es una propiedad de un anillo local
noetheriano $A$ que es una $k$-álgebra; se supone que, si $k'$ es un
cuerpo isomorfo a $k$ y $A'$ es un anillo local noetheriano que es una
$k'$-álgebra di-isomorfa a $A$, las propiedades $\boldsymbol Q(A,k)$
y $\boldsymbol Q(A',k')$ son equivalentes.
\end{env}

\oldpage[IV-2]{193}

Si $X$ e $Y$ son dos preesquemas localmente noetherianos, diremos que
un morfismo $f:X\to Y$ es un \emph{$\boldsymbol P$-morfismo} si:
\begin{enumerate}
\item[1\textsuperscript{o}] $f$ es \emph{plano};
\item[2\textsuperscript{o}] para todo $y\in Y$, se cumple la propiedad
$\boldsymbol P(f^{-1}(y),\boldsymbol{k}(y))$.
\end{enumerate}
